\documentclass[12pt]{article}
\usepackage[utf8]{inputenc}
\usepackage[T1]{fontenc}
\usepackage[margin=1in]{geometry}
\usepackage{setspace}
\usepackage{xcolor}
\usepackage[colorlinks=true, urlcolor=blue!60!black, linkcolor=blue!60!black]{hyperref}
\setlength{\parindent}{2em}
\setlength{\parskip}{1.25em}
\renewcommand{\baselinestretch}{1.0}
\begin{document}
\begin{center}
{\LARGE \textit{Virginia v. Black}, 538 U.S. 343 (2003)\par}
\vspace{0.8em}
{\large Case Brief --- Four-Page Edition\par}
\end{center}
\vspace{1em}

\section{Memory Jogger}

Reviewing Virginia's cross-burning statute, the Court held that a State may ban cross burning carried out with intent to intimidate, a form of true threat, without violating \textit{R.A.V.}; but four Justices would have held the prima facie provision facially invalid as construed in Barry Black's jury instruction. Seven Justices agreed that Black's conviction could not stand (\href{https://supreme.justia.com/cases/federal/us/538/343/}{538 U.S. 343}).

\section{Detailed Case Facts}

Virginia made it a felony ``for any person or persons, with the intent of intimidating any person or group of persons, to burn, or cause to be burned, a cross on the property of another, a highway or other public place.'' A 1968 sentence added: ``Any such burning of a cross shall be prima facie evidence of an intent to intimidate a person or group of persons.''

\textbf{Barry Black} led a Ku Klux Klan rally of 25--30 people on private property, with the owner's permission, visible from a highway; after racist and violent speeches, a 25- to 30-foot cross was burned. Over his objection, the jury was told that ``the burning of a cross by itself is sufficient evidence from which you may infer the required intent,'' and it fined him \$2,500. \textbf{Richard Elliott and Jonathan O'Mara}, unaffiliated with the Klan, tried to burn a cross on the lawn of Elliott's African-American neighbor, James Jubilee, to ``get back'' at him for complaining about gunfire. O'Mara pleaded guilty, reserving his challenge; Elliott's jury received no prima facie instruction.

\section{Procedural History}

The Virginia Court of Appeals affirmed all three convictions. The Supreme Court of Virginia struck down the statute on its face, holding under \textit{R.A.V. v. City of St. Paul} that singling out cross burning was content and viewpoint discrimination and, alternatively, that the prima facie provision made the law overbroad. The U.S. Supreme Court granted certiorari, heard argument on December 11, 2002, and decided the case on April 7, 2003.

\section{Judicial Votes}

Affirmed as to Black (7--2 that his conviction could not stand); vacated and remanded as to Elliott and O'Mara (5--4).
\begin{itemize}
\item \textbf{Majority (Parts I--III):} Justice O'Connor, joined by Chief Justice Rehnquist and Justices Stevens, Scalia, and Breyer.
\item \textbf{Plurality (Parts IV--V):} O'Connor, joined by Rehnquist, Stevens, and Breyer.
\item \textbf{Concurrence:} Justice Stevens.
\item \textbf{In part and dissenting in part:} Justice Scalia (joined by Thomas as to Parts I and II); Justice Souter, joined by Justices Kennedy and Ginsburg.
\item \textbf{Dissent:} Justice Thomas.
\end{itemize}

\section{Holding}

A State may, consistent with the First Amendment and \textit{R.A.V.}, ban cross burning with intent to intimidate, because such intimidation is a ``true threat'' and cross burning is a ``particularly virulent form of intimidation.'' The plurality held the prima facie provision, as construed in Virginia's model instruction, facially unconstitutional because it permits conviction on the act alone; Elliott's and O'Mara's cases were remanded for consideration of severability or a narrowing construction.

\section{Analysis of Opinions}

\textbf{O'Connor (majority and plurality).} Cross burning became ``inextricably intertwined with the history of the Ku Klux Klan,'' serving both as a tool of terror and as a symbol of ideology; when aimed at a target, ``few if any messages are more powerful.'' Though offensive symbolic conduct is protected under \textit{Texas v. Johnson}, true threats are not, and the speaker need not intend to carry out the threat. \textit{R.A.V.} allows content discrimination when its basis ``consists entirely of the very reason the entire class of speech at issue is proscribable''; unlike St.~Paul's ordinance targeting disfavored topics, Virginia banned all intimidating cross burning regardless of motive or victim. For the plurality, the prima facie clause ``strips away the very reason why a State may ban cross burning,'' permitting conviction whenever a defendant presents no defense and ignoring that cross burning may be political ideology at a rally or art in a film.

\textbf{Stevens (concurring).} Intimidating cross burning is unprotected whether or not the State bans other threats.

\textbf{Scalia.} Under Virginia law ``prima facie evidence'' only lets the prosecution survive a motion to strike, so the clause is not substantially overbroad; the plurality wrongly relied on a non-binding instruction. Black's conviction failed, but retrial should be allowed.

\textbf{Souter (with Kennedy and Ginsburg).} Because the burning cross is tied to white-supremacist ideology, the statute does not fit \textit{R.A.V.}'s exception, and the prima facie clause suggests ``official suppression of ideas is afoot.''

\textbf{Thomas (dissenting).} Intimidating cross burning is ``conduct, not expression,'' and the clause is a permissible inference.

\textbf{Later development.} \textit{Counterman v. Colorado} (2023) requires proof of at least recklessness in true-threats prosecutions, reading \textit{Black}'s ``intent'' as Virginia's statutory element.

\section{Examples: Future Applications}

\begin{enumerate}
\item \textbf{Same-Side (Valid): A Noose on an Election Clerk's Porch.} A ban on hanging a noose with intent to intimidate fits the virulence exception, and a conviction can stand if intent is proved without a prima facie instruction.
\item \textbf{Same-Side (Valid): An Effigy Burned in a Driveway.} A general ban on burning objects on another's property to cause fear of bodily harm targets intimidation, not a message.
\item \textbf{Opposite-Side (Invalid): A Noose in a Protest Sculpture.} Convicting an anti-death-penalty artist on the prima facie clause alone raises Black's successful objection.
\item \textbf{Opposite-Side (Invalid): Viewpoint Selection or Objective Standard.} A ban limited to threats based on race, religion, or party may raise \textit{R.A.V.} problems, and a prosecution without proof of recklessness violates \textit{Counterman}.
\item \textbf{Fence-Sitter: Burning-Cross Image Online.} Posting a burning-cross image tagging a Black candidate may be a threat or political expression, and context leaves intent uncertain.
\end{enumerate}

\section{Critique}

\textbf{Tension with \textit{R.A.V.}} As Justice Souter argued, if the burning cross is historically fused with the ideology of white supremacy, upholding a cross-specific ban under the virulence exception may collapse the line between targeting harm and targeting a disfavored viewpoint.

\textbf{Procedural overbreadth.} Scalia's critique raises whether an evidentiary inference in a jury instruction can render an otherwise valid prohibition facially void before a state supreme court has construed it.

\section{Key Quotations}

\begin{quote}
``\,`True threats' encompass those statements where the speaker means to communicate a serious expression of an intent to commit an act of unlawful violence to a particular individual or group of individuals. \dots Intimidation in the constitutionally proscribable sense of the word is a type of true threat, where a speaker directs a threat to a person or group of persons with the intent of placing the victim in fear of bodily harm or death.''
\end{quote}
\noindent --- Justice O'Connor, Opinion of the Court, 538 U.S. at 359--360.

\begin{quote}
``The First Amendment permits Virginia to outlaw cross burnings done with the intent to intimidate because burning a cross is a particularly virulent form of intimidation.''
\end{quote}
\noindent --- Justice O'Connor, Opinion of the Court, 538 U.S. at 363.

\begin{quote}
``The act of burning a cross may mean that a person is engaging in constitutionally proscribable intimidation. But that same act may mean only that the person is engaged in core political speech. \dots The prima facie evidence provision in this case ignores all of the contextual factors that are necessary to decide whether a particular cross burning is intended to intimidate. The First Amendment does not permit such a shortcut.''
\end{quote}
\noindent --- Justice O'Connor, plurality, 538 U.S. at 365, 367.

\end{document}
